\documentclass[12pt]{article}
\usepackage[utf8]{inputenc}
\usepackage[T1]{fontenc}
\usepackage{hyperref}
\title{Southern Katang Language Resources for K-12 Curriculum Development}
\author{Compiled by Claude Sonnet 5.5 (claude-sonnet-5-5) for Andr\'as Kornai}
\date{}
\begin{document}
\maketitle

\section{Introduction}
Southern Katang (ISO 639-3 sct) is a Katuic (Austroasiatic) language of Laos. Glottolog lists it as not endangered, and the Alphabetical Synonym Journal Project (ASJP) estimates about 40,000 speakers. Joshua Project notes that Northern and Southern Katang groups live in Laos and that the Southern group speaks a slightly different language. The searches found no dictionary, primer, textbook or school programme for Southern Katang. This dossier is therefore short and gives only background and comparative sources. Do not confuse this language with the Ta'Oi/Kataang variety or with the Bru dialect called Katang.

\section{Primary Educational Resources}
None found. Only reference and background pages turned up in the searches.

\section{Linguistic and Reference Materials}
\begin{itemize}
\item Glottolog entry for ISO sct. \url{https://glottolog.org/resource/languoid/iso/sct}
\item ASJP wordlist page for Katang Southern. \url{https://asjp.clld.org/languages/KATANG_SOUTHERN}
\item Wikipedia, Katuic languages. \url{https://en.wikipedia.org/wiki/Katuic_languages}
\item Gehrmann's ANU thesis on West Katuic comparative phonology (PL-C132). It contains comparative material and wordlists for Bru, Kui, Pakoh and Katu, and none specifically for Katang. \url{https://openresearch-repository.anu.edu.au/bitstream/1885/145890/1/PL-C132.pdf}
\item Sidwell (2005), \emph{The Katuic Languages: Classification, Reconstruction and Comparative Lexicon}, a bibliographic record. \url{https://glottolog.org/resource/reference/id/134737}
\end{itemize}

\section{Religious and Community Resources}
Joshua Project has a people-group profile for Southern Katang in Laos. \url{https://joshuaproject.net/people_groups/21062/LA}. The Wikipedia article on the Katang people gives ethnographic background. \url{https://en.wikipedia.org/wiki/Katang_people}. No scripture or community organization was found.

\section{Digital and Audio Resources}
None found.

\section{Curriculum Development Context}
Lao is the language of schooling in Laos. No evidence of Southern Katang instruction was found. A curriculum would need fieldwork first: a phonological analysis, a practical orthography (probably Lao-script based) and a lexicon. The Katuic comparative literature above may give a starting point.
\end{document}